% 来源及取材范围见分题文件；此为整理者转录的正文，不是下载原件。
\begin{lemma}[10.81.1]
Let $M$ be an $R$-module. The following are equivalent:
\begin{enumerate}
\item $M$ is flat.
\item If $f:R^n\to M$ is a module map and $x\in\operatorname{Ker}(f)$, then there are module maps $h:R^n\to R^m$ and $g:R^m\to M$ such that $f=g\circ h$ and $x\in\operatorname{Ker}(h)$.
\item Suppose $f:R^n\to M$ is a module map, $N\subset\operatorname{Ker}(f)$ any submodule, and $h:R^n\to R^m$ a map such that $N\subset\operatorname{Ker}(h)$ and $f$ factors through $h$. Then given any $x\in\operatorname{Ker}(f)$ we can find a map $h':R^n\to R^{m'}$ such that $N+Rx\subset\operatorname{Ker}(h')$ and $f$ factors through $h'$.
\item If $f:R^n\to M$ is a module map and $N\subset\operatorname{Ker}(f)$ is a finitely generated submodule, then there are module maps $h:R^n\to R^m$ and $g:R^m\to M$ such that $f=g\circ h$ and $N\subset\operatorname{Ker}(h)$.
\end{enumerate}
\end{lemma}
\begin{proof}
That (1) is equivalent to (2) is just a reformulation of the equational criterion for flatness. To show (2) implies (3), let $g:R^m\to M$ be the map such that $f$ factors as $f=g\circ h$. By (2) find $h'':R^m\to R^{m'}$ such that $h''$ kills $h(x)$ and $g:R^m\to M$ factors through $h''$. Then taking $h'=h''\circ h$ works. (3) implies (4) by induction on the number of generators of $N\subset\operatorname{Ker}(f)$ in (4). Clearly (4) implies (2).
\end{proof}

\begin{lemma}[10.81.2]
Let $M$ be an $R$-module. Then $M$ is flat if and only if the following condition holds: if $P$ is a finitely presented $R$-module and $f:P\to M$ a module map, then there is a free finite $R$-module $F$ and module maps $h:P\to F$ and $g:F\to M$ such that $f=g\circ h$.
\end{lemma}
\begin{proof}
This is just a reformulation of condition (4) from Lemma 10.81.1.
\end{proof}

\begin{theorem}[10.81.4, Lazard's theorem, Tag 058G]
Let $M$ be an $R$-module. Then $M$ is flat if and only if it is the colimit of a directed system of free finite $R$-modules.
\end{theorem}
\begin{proof}
A colimit of a directed system of flat modules is flat, as taking directed colimits is exact and commutes with tensor product. Hence if $M$ is the colimit of a directed system of free finite modules then $M$ is flat.

For the converse, first recall that any module $M$ can be written as the colimit of a directed system of finitely presented modules, in the following way. Choose a surjection $f:R^I\to M$ for some set $I$, and let $K$ be the kernel. Let $E$ be the set of ordered pairs $(J,N)$ where $J$ is a finite subset of $I$ and $N$ is a finitely generated submodule of $R^J\cap K$.

Then $E$ is made into a directed partially ordered set by defining $(J,N)\leq(J',N')$ if and only if $J\subset J'$ and $N\subset N'$. Define $M_e=R^J/N$ for $e=(J,N)$, and define $f_{ee'}:M_e\to M_{e'}$ to be the natural map for $e\leq e'$. Then $(M_e,f_{ee'})$ is a directed system and the natural maps $f_e:M_e\to M$ induce an isomorphism $\operatorname{colim}_{e\in E}M_e\xrightarrow{\cong}M$.

Now suppose $M$ is flat. Let $I=M\times\mathbf Z$, write $(x_i)$ for the canonical basis of $R^I$, and take in the above discussion $f:R^I\to M$ to be the map sending $x_i$ to the projection of $i$ onto $M$. To prove the theorem it suffices to show that the $e\in E$ such that $M_e$ is free form a cofinal subset of $E$. So let $e=(J,N)\in E$ be arbitrary.

By Lemma 10.81.2 there is a free finite module $F$ and maps $h:R^J/N\to F$ and $g:F\to M$ such that the natural map $f_e:R^J/N\to M$ factors as $R^J/N\xrightarrow{h}F\xrightarrow{g}M$. We are going to realize $F$ as $M_{e'}$ for some $e'\geq e$.

Let $\{b_1,\ldots,b_n\}$ be a finite basis of $F$. Choose $n$ distinct elements $i_1,\ldots,i_n\in I$ such that $i_\ell\notin J$ for all $\ell$, and such that the image of $x_{i_\ell}$ under $f:R^I\to M$ equals the image of $b_\ell$ under $g:F\to M$. This is possible since every element of $M$ can be written as $f(x_i)$ for infinitely many distinct $i\in I$ (by our choice of $I$).

Now let $J'=J\cup\{i_1,\ldots,i_n\}$, and define $R^{J'}\to F$ by $x_i\mapsto h(x_i)$ for $i\in J$ and $x_{i_\ell}\mapsto b_\ell$ for $\ell=1,\ldots,n$. Let $N'=\operatorname{Ker}(R^{J'}\to F)$. Observe:
\begin{enumerate}
\item The square
\[\xymatrix{R^{J'}\ar[r]\ar@{^{(}->}[d]&F\ar[d]^g\\R^I\ar[r]_f&M}\]
is commutative, hence $N'\subset K=\operatorname{Ker}(f)$;
\item $R^{J'}\to F$ is a surjection onto a free finite module, hence it splits and so $N'$ is finitely generated;
\item $J\subset J'$ and $N\subset N'$.
\end{enumerate}
By (1) and (2) $e'=(J',N')$ is in $E$, by (3) $e'\geq e$, and by construction $M_{e'}=R^{J'}/N'\cong F$ is free.
\end{proof}
